\section{Methods}\label{sec:methods}

\subsection{MOEA/D with a pool of operators}\label{sec:moea_d_with_a_pool_of_operat}

We use MOEA/D with N weight vectors $\boldsymbol{\lambda}^{i}$ = ($\lambda_{i}$, 1 $-$ $\lambda_{i}$), $\lambda_{i}$ equally spaced in [0.001, 0.999], and the Tchebycheff scalarisation for maximisation,

\begin{equation*}
g(S \mid \boldsymbol{\lambda}, z^{*}) = \max_{k}\; \lambda_k \bigl(z^{*}_k - f_k(S)\bigr),
\end{equation*}

\noindent where $z^{*}$ is the running ideal point. Each subproblem keeps one solution; the T nearest weight vectors form its neighbourhood. At every step a parent is taken from the neighbourhood with probability 0.9 and from the whole population otherwise; one operator is applied; the child replaces every neighbour for which it is not worse ($g_{\mathrm{child}}$ $\leq$ $g_{\mathrm{neighbour}}$). An external archive keeps the non-dominated solutions found (duplicates within 10$^{-9}$ rejected). Subproblem i is initialised by a weight-guided GRASP construction (restricted candidate list with $\alpha$ = 0.3 on $\lambda$$_{1}$$\Delta$$_{1}$ + $\lambda$$_{2}$$\Delta$$_{2}$, where $\Delta$$_{1}$ is the candidate's summed distance to the current set and $\Delta$$_{2}$ its minimum distance, in raw units). The population size N, neighbourhood size T and number of generations depend on n (Appendix A).

Algorithm 1 summarises one run (a generation is N steps, one per subproblem).

\begin{algorithm}[t]
\caption{MOEA/D with a pool of operators and a selector}
\label{alg:moead}
\begin{algorithmic}
\REQUIRE instance, $N$ weight vectors, neighbourhood size $T$, $G$ generations, pool of 8 operators, selector $\sigma$
\STATE for $i=1,\dots,N$: build $S_i$ by weight-guided GRASP and evaluate it
\STATE $A \leftarrow$ non-dominated set of $\{S_i\}$; $z^{*}\leftarrow$ component-wise maximum and $z_{\min}\leftarrow$ component-wise minimum of the objectives
\FOR{$g=1,\dots,G$}
\FOR{$i$ in a random permutation of $1,\dots,N$}
\STATE $k \leftarrow \sigma.\mathrm{select}()$ \COMMENT{uniform, UCB1, Thompson sampling or Gaussian sampler}
\STATE $j \leftarrow$ a random member of the neighbourhood of $i$ (probability 0.9) or of the population (0.1)
\STATE $S' \leftarrow \mathrm{operator}_k(S_j)$ (SubsetCrossover also takes a second parent from the neighbourhood); $S' \leftarrow \mathrm{repair}(S')$ if infeasible
\STATE $r \leftarrow$ reward of $S'$ (HV gain with respect to $A$, or Tchebycheff improvement over $S_j$); $\sigma.\mathrm{update}(k, r)$
\FORALL{$l$ in the neighbourhood of $i$}
\IF{$g(S' \mid \lambda_l, z^{*}) \leq g(S_l \mid \lambda_l, z^{*})$}
\STATE $S_l \leftarrow S'$
\ENDIF
\ENDFOR
\STATE update $z^{*}$ and $z_{\min}$; $A \leftarrow$ non-dominated$(A \cup \{S'\})$
\ENDFOR
\ENDFOR
\STATE optional: add the MaxMin anchor solution(s) of Section~\ref{sec:anchor} to $A$
\STATE return $A$
\end{algorithmic}
\end{algorithm}


\textbf{Objective normalisation.} MaxSum and MaxMin live on very different scales ($\approx$10$^{4}$--10$^{6}$ against $\approx$10$^{0}$--10$^{2}$). With raw objectives the $\lambda$$_{1}$-term of g dominates for almost every weight vector, even $\lambda$ = (0.001, 0.999), so the subproblems collapse onto the MaxSum end. In the normalised regime each objective is divided by its running range ($z^{*}_k$ $-$ $z^{\min}_k$) inside g and inside every operator that uses g. All reported results use this regime unless stated otherwise; Section 6.1 quantifies the difference against the raw regime used by our earlier catalogue.

\subsection{The operator pool}\label{sec:the_operator_pool}

The pool has eight operators; after any of them a feasibility repair restores the constraints (Section 4.4):

\begin{enumerate}
\item \emph{RandomSwap1$\times$1} --- remove a random element, insert a random outside element.
\item \emph{BestSwap1$\times$1} --- remove a random element, insert the outside element that maximises MaxSum.
\item \emph{RandomSwap2$\times$2} --- remove two random elements, insert two random outside elements.
\item \emph{GraspGreedy} --- rebuild the solution from scratch with a GRASP construction on the MaxSum contribution, with $\alpha$ drawn uniformly in [0, 0.3] at every call.
\item \emph{BiObjBestSwap} --- remove a random element and insert the outside element that minimises the Tchebycheff value of the child for the weight vector and ideal point of the current subproblem (all outside candidates are evaluated).
\item \emph{MaxMinSwap} --- remove the element that causes the smallest internal distance (the MaxMin bottleneck) and insert the outside element with the best contribution to MaxMin.
\item \emph{LargePerturbation} --- remove k random elements, k uniform in \{3, 4, 5\}, and reinsert them with a GRASP construction ($\alpha$ = 0.2).
\item \emph{SubsetCrossover} --- keep the elements common to the current solution and a partner drawn from the neighbourhood, shuffle the elements exclusive to either parent and add them until the cardinality m (MDP) or the capacity requirement (CDP, GDP; within the budget) is met.
\end{enumerate}

\subsection{Operator selection strategies}\label{sec:operator_selection_strategie}

\emph{Uniform} selection draws one of the eight operators uniformly at random; it has no memory and no parameters. The adaptive strategies are UCB1 \citep{auer2002} with a sliding window (window 50, exploration constant c = 0.5), Thompson sampling \citep{thompson1933} with Beta posteriors over a sliding window of binary successes, and a Gaussian sliding-window Thompson sampler (GTS) with a normal-inverse-$\chi$$^{2}$ posterior over continuous rewards. (The sampler appears as \texttt{dyts\_\allowbreak{}*} in our code and result files because it was first meant to follow the dynamic Thompson sampling of \citet{sun2020}; it is \emph{not} that algorithm, which uses Beta posteriors, a binary reward and a decay threshold C. The original is implemented separately and tested in Section 6.4.) Two historical reward definitions are used: \emph{HV} --- the hypervolume gain of the child with respect to the archive (zero if the child is dominated), and \emph{Scal} --- the improvement of the child over its parent in the subproblem's Tchebycheff value. This gives six adaptive MOEA/D variants plus Uniform, i.e. seven MOEA/D configurations. Section 6.4 adds two further reward signals, cost-aware scaling and an exploration floor.

\emph{Success criteria and floating-point tolerance.} The Thompson sampler and the DYTS count a child as a success when its reward is positive, and the insert reward of Section 6.4 gives 1 when no archive member weakly dominates the child. The same set of elements evaluated in a different order gives objective values that differ by about 10$^{-}$$^{1}$$^{5}$ in relative terms, which happens whenever a child reproduces a member of the archive or of the neighbourhood (typically after SubsetCrossover). We therefore require a gain larger than 10$^{-9}$ in relative terms for a success, and treat a point within that tolerance of an archive member as blocked. Our first runs lacked this tolerance: the Thompson variants then spent 0.20--0.24 of their selections on SubsetCrossover (uniform selection: 0.125), mostly on spurious successes. All results of the variants that depend on the criteria --- the two Thompson variants of the catalogue, the insert reward, the published DYTS and its binary variants, FRRMAB and the hyper-parameter and contextual variants built on them --- were recomputed with the tolerance; the UCB1 and Gaussian-sampler variants, whose rewards are continuous, are unaffected (their runs are identical bit for bit).

\subsection{Feasibility repair (CDP, GDP)}\label{sec:feasibility_repair_cdp_gdp}

After any operator, if the solution is infeasible it is repaired in two phases: elements are removed in increasing order of efficiency (capacity per unit cost; capacity alone in CDP) until the budget holds (GDP only), then outside elements are added in decreasing order of efficiency, skipping those that would break the budget, until the capacity requirement is met. A fallback constructive repair guarantees feasibility. The repair is applied identically to all methods that need it.

\subsection{MaxMin anchoring}\label{sec:anchor}

To cover the MaxMin end of the front we add an optional final step. A single-objective tabu search starts from the archive's best-MaxMin solution and from farthest-point constructions; each iteration removes one element of a bottleneck pair, rebuilds the solution by repeatedly adding the element farthest from the current set (random choice among the best three with probability 0.3; recently removed elements are tabu for 10--19 iterations), and prunes redundant elements --- removing an element never decreases MaxMin, so good solutions are capacity-minimal. In CDP and GDP rebuilding respects the budget. We use 2,000 iterations $\times$ 3 restarts (6,000 $\times$ 3 for n = 500) and insert the result into the archive. Cost: $\approx$2 s per run for n $\leq$ 150.

\emph{Solution-based variant.} Because the published best values on the hardest family (MDG-b) come from tabu searches over whole solutions, we also implemented a solution-based tabu search in the style of \citet{lu2023} (our own NumPy reimplementation from the paper's description, not the authors' code). From a greedy farthest-point solution it moves at every iteration to the best non-tabu neighbour in the union of three neighbourhoods --- add an element, drop an element that keeps the capacity, and a constrained swap that exchanges an element of the bottleneck pair with one of the eight outside elements farthest from the rest of the solution among those that keep capacity (and, in the GDP, the budget) feasible --- and a solution is tabu if it was visited before (exact hash set; if every neighbour is tabu the best one is taken). We use 3,000 iterations for n $\leq$ 150 and 30,000 for n = 500, either alone (\texttt{SBTS}; three restarts at n = 500) or as a second step after the tabu anchor, started from the archive's best-MaxMin solution (\texttt{combined}; tabu 2,000 $\times$ 3 or 6,000 $\times$ 3 iterations, then one SBTS run of 3,000 or 30,000 iterations; both results enter the archive).

\subsection{Baselines}\label{sec:baselines}

\emph{NSGA-II} \citep{deb2002}: tournament selection on rank and crowding distance, subset crossover (probability 0.9, same recombination as operator 8), swap mutation (probability 0.1), GRASP initialisation, feasibility repair, and an external archive. \emph{SMS-EMOA} enters only in Section 4.7. \emph{GRASP-MO}: at each iteration a weight vector is drawn from a Dirichlet(1,1) distribution, a solution is built by a GRASP construction \citep{feo1995} guided by the weighted sum ($\alpha$ = 0.3, respecting capacity and budget), improved by one best-improvement 1$\times$1 swap pass on the same scalarisation, and inserted into a non-dominated archive. \emph{GRASP-PR} adds a final path-relinking phase over five random archive pairs. \emph{VNS-Pareto} is a variable neighbourhood search \citep{hansen2001} over three neighbourhoods of increasing size (a 1-swap, a 2-swap and a 6-element destroy-and-greedy-rebuild) started from archive members, with early stopping after max(50, 0.3 $\times$ iterations) iterations without archive change; \emph{VNS-Tchebycheff} runs the same neighbourhoods per weight vector on the Tchebycheff value. The same families of methods (NSGA-II, GRASP, VNS) were used for the bi-objective MDP by \citet{colmenar2018}, but we know of no published bi-objective MaxSum + MaxMin implementation for CDP/GDP, and we did not have access to their code, so all baselines are our own adaptations, with budgets fixed by the rule of Appendix A and no parameter tuning in the main experiment; the NSGA-II configuration that Colmenar et al. obtained with irace (population 340, 28,130 generations, about 10$^{7}$ individuals produced per run) is roughly 500 times larger than ours (80 $\times$ 250 at n = 500), and their best method, a tabu search that moves between non-dominated solutions with a hash-based tabu list, enters our comparison in Section 4.7; we use the same instance, repair, normalisation (where applicable) and archive code for all methods.

\subsection{Trajectory baselines of Colmenar et al. and the tuning protocol}\label{sec:trajectory_baselines_of_colm}

The baselines of Section 4.6 are the ones used in the main experiment, with fixed budgets and no tuning. Because \citet{colmenar2018} tuned every method with irace \citep{lopezibanez2016} and found a tabu search to be the best one, we added two trajectory methods in their scheme (\texttt{bo\_\allowbreak{}trajectory.py}; our reimplementation from the description in their paper, not their code), SMS-EMOA, and a tuning protocol applied to \emph{all} methods, MOEA/D included.

\emph{Tabu search (TS) and VNS with the front as incumbent.} Following the scheme of \citet{duarte2015} that Colmenar et al. use, the approximation of the Pareto front is the incumbent solution: an iteration explores the neighbourhood of the current solution S with a local search that feeds the archive; if the archive received no insertion there was no improvement, S becomes a random non-tabu member of the archive and the stagnation counter grows. The constructive is a GRASP with restricted candidate list (g $\geq$ $g_{\mathrm{min}}$ + $\alpha$($g_{\mathrm{max}}$ $-$ $g_{\mathrm{min}}$)) that alternates the objective between constructions (\texttt{altbws}) or at every inserted element (\texttt{altins}), or is random. The local search (\texttt{pr}) hill-climbs MaxMin from S and MaxSum from S with first-improvement swaps (evaluated in batch per removed element, lexicographic in the other objective), then path-relinks the two results with random steps; every traversed solution is offered to the archive. Alternatives are a local search that alternates the objective per call (\texttt{alt}) and one that accepts only dominating swaps (\texttt{dom}). In the TS a solution is tabu for \texttt{tenure} iterations after it is visited or selected (exact Zobrist hash); in the VNS the shake exchanges k random elements, k grows to $k_{\mathrm{max}}$ and returns to 1 when the archive changes. With a time limit the original stopping rule becomes a restart: after \texttt{max\_\allowbreak{}no\_\allowbreak{}improve} stagnant iterations (TS) or \texttt{restart\_\allowbreak{}after} complete cycles of k (VNS) a new solution is constructed. On the CDP and GDP, swaps that violate capacity or budget are discarded, the MaxMin local search also accepts removing an element when the capacity is still met and the MaxMin improves strictly, redundant elements are pruned at the end (both points enter the archive) and path-relinking repairs infeasible intermediates. The ``irace setting'' rows use the values of Table 6 of Colmenar et al. (TS: $\alpha$ = 0.94, constructive \texttt{altbws}, local search \texttt{alt}, tenure 64\% of the number of elements --- our reading of ``tenure configured as a percentage'' ---, 70 iterations without improvement; VNS: $\alpha$ = 0.19, \texttt{altins}, \texttt{pr}, $k_{\mathrm{max}}$ = 29.49\% of the solution size, seven cycles of k before a restart).

\emph{SMS-EMOA.} For two objectives the steady-state, hypervolume-based SMS-EMOA \citep{beume2007} is the natural alternative to NSGA-II. At every step it produces one child with the NSGA-II operators (random parents, subset crossover, swap mutation, repair, GRASP initialisation), adds it to the population and removes, from the last non-dominated front, the individual with the smallest hypervolume contribution (computed on objectives scaled by the range of the population, with the reference point 0.1 below the minimum); the external archive is the same as in NSGA-II.

\emph{Tuning.} Each method is tuned separately for each of six cells (CDP, GDP and MDP; n $\leq$ 150 and n = 500) by random search with successive halving on eight training instances that are excluded from the test set (seven for the CDP with n = 500): 16 candidate configurations (the default ones plus random draws from the space below); the first stage runs all candidates on two training instances, the second the best six on three more, the third the best two on the last three; the winner has the highest mean $\mathrm{HV}_{\mathrm{norm}}$ over the instances it ran on. Every evaluation runs with the wall-clock limit equal to the time the default MOEA/D-Uniform needs on that instance. The spaces are: NSGA-II and SMS-EMOA --- population in [20, 400] (log scale), crossover probability in [0.2, 1], mutation probability in [0.02, 0.9] (log), initialisation GRASP or random, GRASP $\alpha$ in [0.05, 0.5]; GRASP --- $\alpha$ in [0, 0.8], 0--3 local-search passes, path relinking on or off, 2--20 relinked pairs; TS --- $\alpha$ in [0.3, 1], constructive, local search, tenure in [5, 200], stagnation limit in [5, 300]; VNS --- $\alpha$ in [0, 1], constructive, local search, $k_{\mathrm{max}}$ in [2, 40], restart limit in [1, 20]; MOEA/D-Uniform --- number of weight vectors in [10, 200], neighbourhood size 5--40\% of that number, GRASP $\alpha$ of the initialisation in [0.05, 0.8], with the pool of eight operators and uniform selection fixed. The tuning uses about 56 runs per method and cell (336 over the six cells) and is the same for every method; Colmenar et al. capped irace at 500 experiments per algorithm on 29 instances and report that more experiments gave similar results. The default MOEA/D, the one of the main experiment, is also reported.

\subsection{Size-aware operators, a cardinality bound and trajectory seeding}\label{sec:size_aware_operators_a_cardi}

Three additions, all opt-in (the standard pool and the standard problem are unchanged and reproduce the stored runs bit for bit), address findings of Sections 6.12 and 6.13.

\emph{Size-aware operators (CDP, GDP).} None of the eight operators of Section 4.2 chooses the size of the set: the swaps keep it, the constructions stop at the capacity requirement, and the repair adds elements only when a move breaks the requirement. Two operators that do are added to the pool (Section 6.14). BiObjGrow inserts the outside element that minimises the Tchebycheff value of the current subproblem, evaluating all candidates in one batch (the budget on the GDP and the cardinality bound below are respected, and the operator leaves at least three outside elements so that the swap operators stay applicable). BottleneckShrink removes one endpoint of the closest pair of the set, the one whose removal leaves the larger MaxMin, provided that the capacity requirement still holds; otherwise it falls back to MaxMinSwap.

\emph{Cardinality-bounded variant (CDP, GDP).} Because f$_{1}$ grows with the size of the set (Section 3), the bound |S| $\leq$ $\lceil$$\rho$ $k_{\mathrm{ref}}$$\rceil$ makes the bi-objective problem well posed, with a finite trade-off between small and larger sets. $k_{\mathrm{ref}}$ is the size of the solution returned by the efficiency-greedy construction (elements ordered by capacity per unit of cost), which is feasible by construction on the GDP as well, and $\rho$ = 2 ($\rho$ = 1.5 in a sensitivity run). Every method obeys the bound: constructions and the repair remove, one at a time, the member with the smallest contribution to MaxSum among those that can leave without breaking the capacity requirement (falling back to the greedy solution if none can), and the dominating insertions of tabu search and VNS stop at the bound. The instance classes expose it as \texttt{set\_\allowbreak{}max\_\allowbreak{}size($\rho$)}; with the default (\texttt{None}) nothing changes.

\emph{Trajectory seeding.} In the hybrid variants (E20, Section 6.15) a tabu search or VNS with the tuned parameters of Section 4.7 runs for a fraction of the time limit and its archive becomes the initial population of the MOEA/D (each archive member replaces the subproblems for which it has a smaller Tchebycheff value); the MOEA/D then runs for the rest of the time and the result is the non-dominated union. The seeding phase counts in the time limit, so every method has the same wall-clock budget. A scalarised local-search operator (hill-climbing on the Tchebycheff value of the subproblem until a local optimum or ten swaps) was also tried as a ninth pool member.
